\section{Representation alignment: do independent models share an algebra?}\label{sec:align}
\paragraph{Procedure.}
All models of a group share the observation map, so we can encode a common probe set (200 fresh sequences) with
every model. For a seed pair $(a,b)$ we fit the orthogonal Procrustes map $R$ from $E_a(x)$ to $E_b(x)$ on that
probe set, averaging both directions. We then measure four quantities:
\begin{itemize}
\item latent linear CKA;
\item the \emph{cross-seed algebra distance} between $\spn(RA^{(a)}R^\top)$ and $\spn(A^{(b)})$;
\item the disagreement of inferred actions, $\|R\rho_aR^\top-\rho_b\|_F^2$, normalized by $\|\rho_b-I\|_F^2$;
\item \emph{transformation transfer}: on a disjoint held-out probe set, the action inferred by $a$ is applied in
$b$'s latent space, $\hat x_{t+1}=D_b(R\rho_a(x_t,x_{t+1})R^\top E_b(x_t))$. We report the transfer gap
$(\mathrm{MSE}_{\text{transfer}}-\mathrm{MSE}_{\text{native}})/(\mathrm{MSE}_{\text{identity}}-\mathrm{MSE}_{\text{native}})$,
which is 0 when transfer is as good as $b$'s own action and 1 when it is no better than doing nothing.
\end{itemize}
Units are the \NNpairs{} disjoint seed pairs, identical for every method.

\begin{table}[t]
\centering\scriptsize
\setlength{\tabcolsep}{3.5pt}
\caption{Cross-seed alignment on \NNpairs{} disjoint seed pairs (mean$\pm$s.d.). Algebra distance: chance
$0.80$ ($T^2$), $0.50$ (SO(3)). Transfer gap: 0 = native quality, 1 = identity guess. ``Projected'' re-expresses
the transferred action in $b$'s generator basis before exponentiating.}
\label{tab:align}
\begin{tabular}{llccccc}
\toprule
Group & Method & CKA & Algebra dist. & Action disagr. & Transfer gap & Projected gap\\
\midrule
\tabinput{../../results/final/tables/alignment.tex}
\bottomrule
\end{tabular}
\end{table}

\paragraph{Algebras agree more across seeds, and CKA does not see it.}
For SO(3), the cross-seed algebra distance falls from \NHthreeSOthreeMeanb{} (+Comp.) to \NHthreeSOthreeMeana{}
(BracketNet), lower in \NHthreeSOthreeWins{} pairs (Holm $p=\NHthreeSOthreeHolmp$; Table~\ref{tab:primary}). For
$T^2$ the difference is not significant (Holm $p=\NHthreeTtwoHolmp$). Over the same pairs, latent CKA is
\NAlSOthreeCompCka{} vs.\ \NAlSOthreeBNCka{} for SO(3) and \NAlTtwoCompCka{} vs.\ \NAlTtwoBNCka{} for $T^2$
(Table~\ref{tab:align}). Pointwise similarity of embeddings carries essentially no information about whether two
models represent transformations by the same algebra.

\paragraph{Why agreement improves: disagreement becomes discrete.}
Classifying each BracketNet SO(3) pair by the categories of its two models (Fig.~\ref{fig:align}) gives:
\begin{itemize}
\item \NPtSOthreeBNBothcorrect{} pair with both models correct;
\item \NPtSOthreeBNBothsameincorrecttype{} pair with both models in the \emph{same} incorrect (chiral) class, at
distance $0$: agreement on a wrong algebra;
\item \NPtSOthreeBNCloseddifferenttypes{} pairs that pair a correct model with a chiral one;
\item \NPtSOthreeBNInvolvesnonclosedorcollapsed{} pairs involving a non-closed model.
\end{itemize}
The correct/chiral pairs sit at distance $\approx0.5$. This value is fixed by the geometry: the diagonal
$\so(3)$ and a chiral factor meet at principal angles of $45^\circ$. For \textsc{+Comp.}, all
\NPtSOthreeCompInvolvesnonclosedorcollapsed{} pairs involve a non-closed model, and their distances spread over
$[0,1]$. The reduction in mean distance is therefore real, but it reflects convergence onto a small set of closed
conjugacy classes, as Proposition~\ref{prop:classify} predicts, more than convergence onto the correct algebra.
This is the sense in which closure makes learned structure \emph{comparable} without making it \emph{identifiable}.

\begin{figure}[t]\centering
\includegraphics[width=\linewidth]{fig_alignment.pdf}
\caption{Cross-seed algebra distance on 10 disjoint seed pairs (BN = BracketNet); grey lines connect the same seed pair across methods.
For SO(3), BracketNet pairs concentrate at $0$ (same class) and $\approx0.5$ (diagonal vs.\ chiral), while baseline
pairs spread out. Right: latent CKA is nearly identical for all methods.}\label{fig:align}
\end{figure}

\paragraph{Transfer quality tracks structural correctness.}
Mean transfer gaps barely differ between methods. For BracketNet minus \textsc{+Comp.}, the 95\% intervals are
$[\NApTtwoBNTransfermseonegapCilo,\NApTtwoBNTransfermseonegapCihi]$ for $T^2$ and
$[\NApSOthreeBNTransfermseonegapCilo,\NApSOthreeBNTransfermseonegapCihi]$ for SO(3). What matters is structure.
Pooling the pairs of all methods (a descriptive analysis: the methods share seeds), $T^2$ pairs in which both models
are correct-closed transfer almost perfectly: median gap \NTcTtwoBothcorrectMedian{} over \NTcTtwoBothcorrectN{}
pairs. All other pairs have median gap \NTcTtwoNotbothcorrectMedian{} over \NTcTtwoNotbothcorrectN{} pairs. For
SO(3) only \NTcSOthreeBothcorrectN{} pair is correct on both sides (gap \NTcSOthreeBothcorrectMedian); the remaining
\NTcSOthreeNotbothcorrectN{} have median gap \NTcSOthreeNotbothcorrectMedian. Cross-model reuse of transformations is
thus governed by whether both models recovered the right algebra, which no method here guarantees. The one
near-exception is instructive: an SO(3) +Comp.\ pair that misses the pre-registered closure threshold, but whose
algebras are within $0.01$ of the truth, transfers well (gap $0.03$).
